% ------------------------------------------------------------------------------------
%                           PROBABILITA' E PROCESSI ALEATORI
% ------------------------------------------------------------------------------------
\chapter{Ripasso probabilità e processi aleatori}

Abbiamo fin'ora trattato lo \textbf{studio dei segnali} e dei \textbf{sistemi}
da un \textbf{punto di vista deterministico}, che rappresenta tuttavia
una \textbf{semplificazione}: nella \textbf{realtà} i segnali
sono \textbf{strettamenti legati} alle \textbf{variabili aleatorie} ed
ai \textbf{processi aleatori}.

Questi strumenti ci permettono infatti di \textbf{modellare i disturbi}
(e.g. \textit{interferenze}, \textit{rumore}), che \textbf{non possono} 
essere \textbf{descritti deterministicamente}, rendendo quindi
l'\textbf{intero segnale} assimilabile ad un \textbf{processo aleatorio}.
Questo concetto è mostrato in figura~\ref{fig:07_01_Rumore}.

\begin{figure}[!ht]
    \centering
    % --- Primo Plot: Segnale Pulito ---
    \begin{minipage}{0.32\textwidth}
        \centering
        \begin{tikzpicture}
            \begin{axis}[
                width=\textwidth,
                height=4.5cm,
                axis lines=middle,
                xlabel={$t$},
                ylabel={$s(t)$},
                xmin=0, xmax=2,
                ymin=-1.5, ymax=1.5,
                xtick={0,1,2},
                ytick={-1,0,1},
                grid=major,
                grid style={dashed, gray!30},
                title={Segnale Pulito},
                title style={font=\bfseries\small},
                xlabel style={below right},
                ylabel style={above left},
                ]
                % Segnale sinusoidale ideale
                \addplot[thick, nord10, domain=0:2, samples=200] {sin(deg(2*pi*x))};
            \end{axis}
        \end{tikzpicture}
    \end{minipage}\hfill
    % --- Secondo Plot: Rumore ---
    \begin{minipage}{0.32\textwidth}
        \centering
        \begin{tikzpicture}
            \begin{axis}[
                width=\textwidth,
                height=4.5cm,
                axis lines=middle,
                xlabel={$t$},
                ylabel={$n(t)$},
                xmin=0, xmax=2,
                ymin=-1.5, ymax=1.5,
                xtick={0,1,2},
                ytick={-1,0,1},
                grid=major,
                grid style={dashed, gray!30},
                title={Rumore},
                title style={font=\bfseries\small},
                xlabel style={below right},
                ylabel style={above left},
                ]
                % Rumore simulato (somma di sinusoidi ad alta frequenza)
                \addplot[thick, nord11, domain=0:2, samples=150] 
                    {0.3*sin(deg(2*pi*x*17)) + 0.2*sin(deg(2*pi*x*31)) + 0.15*sin(deg(2*pi*x*47))};
            \end{axis}
        \end{tikzpicture}
    \end{minipage}\hfill
    % --- Terzo Plot: Segnale + Rumore ---
    \begin{minipage}{0.32\textwidth}
        \centering
        \begin{tikzpicture}
            \begin{axis}[
                width=\textwidth,
                height=4.5cm,
                axis lines=middle,
                xlabel={$t$},
                ylabel={$y(t)$},
                xmin=0, xmax=2,
                ymin=-1.5, ymax=1.5,
                xtick={0,1,2},
                ytick={-1,0,1},
                grid=major,
                grid style={dashed, gray!30},
                title={Segnale + Rumore},
                title style={font=\bfseries\small},
                xlabel style={below right},
                ylabel style={above left},
                ]
                % Somma del segnale pulito e del rumore
                \addplot[thick, nord15, domain=0:2, samples=200] 
                    {sin(deg(2*pi*x)) + 0.3*sin(deg(2*pi*x*17)) + 0.2*sin(deg(2*pi*x*31)) 
                    + 0.15*sin(deg(2*pi*x*47))};
            \end{axis}
        \end{tikzpicture}
    \end{minipage}
    \caption{Distorsione di un segnale a causa del rumore}
    \label{fig:07_01_Rumore}
\end{figure}

Un classico esempio è dato dal \textbf{rumore termico}:

\begin{remark}[RUMORE TERMICO]
    Il \textbf{rumore termico} è un esempio di rumore, dovuto
    all'\textit{agitazione termica} delle \textit{particelle}
    nel \textit{mezzo di trasmissione} del segnale.
\end{remark}

Vale allora la pena fare un richiamo veloce dei concetti fondamentali
della probabilità, delle variabili e dei processi aleatori.

\section{Variabili aleatorie}

Ricordiamo la definizione di \textbf{variabile aleatoria}:

\begin{definition}[VARIABILE ALEATORIA]
    Una \textbf{variabile aleatoria} $X$ è una \textbf{funzione}
    \[
        X: \Omega \to E
    \] 
    Dove $\Omega $ è uno \textbf{spazio campionario}.
\end{definition}

Dove per \textbf{spazio campionario} intendiamo 
un \textbf{insieme di possibili eventi}. 

Ricordiamo anche la definizione di \textbf{evento} 
e di \textbf{probabilità di un evento}:

\begin{definition}[EVENTI E LORO PROBABILITA']
    Si dice \textbf{evento} $\omega $ un 
    \textbf{sottoinsieme dello spazio campionario} $\Omega $:
    \[
        \omega \subseteq \Omega 
    \]
    E indichiamo con $P(\omega)$ la \textbf{funzione}:
    \[
        \begin{aligned}
            P: &\Omega \to [0,1] \\ 
               &\omega \mapsto P(\omega ) \in [0,1]
        \end{aligned}
    \]
    Deve \textbf{sempre valere} $P(\Omega) = 1$
\end{definition}

Ricordiamo che $P$ è la \textbf{probabilità} che si 
verifichi l'evento $\omega$, ed è anche chiamata funzione
di \textbf{densità di probabilità}, visto che ci dice
\textit{come la probabilità degli eventi è distribuita nello spazio campionario}.

\begin{remark}
    Se $\omega $ è un \textbf{singolo elemento} di $\Omega $, si
    parla di \textbf{evento elementare}.
\end{remark}

Ricordiamo inoltre che le variabili aleatorie possono essere:
\begin{itemize}
    \item \textbf{CONTINUE} se $\Omega  = [a,b] \subseteq \mathbb{R}$.
    \item \textbf{DISCRETE} se $\Omega  \subseteq \mathbb{Z}$.
\end{itemize}

\subsection{V.A. Continue}

Data una \textbf{v.a. continua} $X$ definita su uno \textit{spazio campionario} $\Omega $, 
la sua \textbf{densità di probabilità} è spesso indicata con $p_X$.
Per la condizione di normalizzazione, deve valere inoltre:
\[
    \int_{x\in \Omega} p_X(x) dx = 1
\]

Si definisce inoltre la \textbf{distribuzione cumulativa} (o \textbf{riparitizione}) 
di $X$ come:

\begin{definition}[FUNZIONE DI DISTRIBUZIONE CUMULATIVA]
    \[
        F_x (u) \coloneqq  P[X \le u] = \int_{-\infty}^u p_X(x) dx
    \] 
\end{definition}

Per come è definita, $F_x$ è \textbf{non decrescente} e abbiamo:
\[
    \begin{aligned}
        \lim_{u \to -\infty} F_x(u) &= 0 \\ 
        \lim_{u \to +\infty} F_x(u) &= 1 \\ 
    \end{aligned}
\]

E' inoltre utile ricordare che valgono:

\begin{prop}
    \[
        \begin{aligned}
            P[x \in [a,b]] &= F_x(b) - F_x(a) \\ 
            p_X(x) &= \frac{d}{dx} F_x (x)
        \end{aligned}
    \] 
\end{prop}

\subsection{V.A. Discrete}

Per quanto riguarda le \textbf{v.a. discrete}, le 
\textbf{definizioni sono analoghe}, con la sola differenza
che gli \textit{integrali sono sostituiti da sommatorie}.

Data una v.a. discreta $X$ definita su $A$, la condizione 
di normalizzazione diventa quindi:
\[
    \sum_{a\in A} p_X(a) = 1
\]
Per la \textbf{distribuzione cumulativa} abbiamo quindi:

\begin{definition}[DISTRIBUZIONE CUMULATIVA]
    \[
        F_x(u) \coloneqq P[X\le u] = \sum_{a\le u, a\in A} p_X(a)
    \] 
\end{definition}

\subsection{Media, varianza e momenti}

Ricordiamo inoltre che, data una v.a. $X$, possiamo definirne
la \textbf{media} come segue:

\begin{definition}[MEDIA]
    La \textbf{media} (o \textbf{valore atteso}) di una v.a. $X$ è:
    \[
        \begin{aligned}
            m_x = m_x(1) = \mathbb{E}[X] &\coloneqq \int_{-\infty}^{\infty} x p_X (x) dx \\ 
            m_x = m_x(1) = \mathbb{E}[X] &\coloneqq \sum_{a_i \in A} a_i p_X(a_i)
        \end{aligned}
    \] 
    Rispettivamente se $X$ è \textbf{continua} oppure \textbf{discreta}.
\end{definition}

Si definisce anche il \textbf{momento secondo} di $X$, ovvero il
\textit{valore atteso del quadrato} di $X$, come segue:

\begin{definition}[MOMENTO SECONDO]
    Il \textbf{momento secondo} di una v.a. $X$ è:
    \[
        \begin{aligned}
            M_x = m_x(2) = \mathbb{E}[X^2] &\coloneqq \int_{-\infty}^{\infty} x^2 p_X (x) dx \\ 
            M_x = m_x(2) = \mathbb{E}[X^2] &\coloneqq \sum_{a_i \in A} a_i^2 p_X(a_i)
        \end{aligned}
    \] 
    Rispettivamente se $X$ è \textbf{continua} oppure \textbf{discreta}.
\end{definition}

In generale, abbiamo:
\[
    m_x(n) = \mathbb{E}[X^n]
\]

Infine, definiamo la \textbf{varianza} di una v.a. come segue:

\begin{definition}[VARIANZA]
    La \textbf{varianza} di una v.a. $X$ è:
    \[
        S_{xx} = \sigma_x ^2 \coloneqq\mathbb{E}[(X - m_x)^2] = \mathbb{E}[X^2] - \mathbb{E}^2[X] \\ 
    \] 
    Rispettivamente se $X$ è \textbf{continua} oppure \textbf{discreta}.
\end{definition}

\subsection{Variabili uniformi e gaussiane}

Le principali \textbf{distribuzioni di probabilità} con cui avremo a che fare 
sono la \textbf{distribuzione uniforme} e la \textbf{distribuzione gaussiana}.

\begin{definition}[DISTRIBUZIONE UNIFORME]
    Diciamo che una v.a. $X$ è \textbf{uniforme} se 
    la sua \textbf{densità} è data da:
    \[
        p_X(x) = \begin{cases}
            \frac{1}{b-a} & a < x < b \\ 
            0 & \text{altrove}
        \end{cases}
    \] 
\end{definition}

Inoltre scriviamo $x\sim \mathcal{U}(a,b)$ per indicare che 
$x$ è campionata da una distribuzione uniforme su $[a,b]$.

\begin{figure}[!ht]
    \centering
    \begin{tikzpicture}
        \begin{axis}[
            width=0.8\textwidth,
            height=6cm,
            axis lines=middle,
            xlabel={$x$},
            ylabel={$p_X(x)$},
            xmin=-4, xmax=4,
            ymin=0, ymax=1.2,
            xtick={-4,-3,-2,-1,0,1,2,3,4},
            ytick={0, 0.2, 0.4, 0.6, 0.8, 1.0},
            grid=major,
            grid style={dashed, gray!30},
            xlabel style={below right},
            ylabel style={above left},
            legend pos=north east,
            legend style={font=\small, draw=none, fill=none},
            ]
            % U(-0.5, 0.5): larghezza 1, altezza 1
            \addplot[thick, nord10, const plot, sharp plot] 
                coordinates {(-4,0) (-0.5,0) (-0.5,1) (0.5,1) (0.5,0) (4,0)};
                \addlegendentry{\(\mathcal{U}(-0.5, 0.5)\)}

            % U(-1.5, 1.5): larghezza 3, altezza 1/3
            \addplot[thick, nord11, const plot, sharp plot] 
                coordinates {(-4,0) (-1.5,0) (-1.5,0.333) (1.5,0.333) (1.5,0) (4,0)};
                \addlegendentry{\(\mathcal{U}(-1.5, 1.5)\)}

            % U(-3, 3): larghezza 6, altezza 1/6
            \addplot[thick, nord14, const plot, sharp plot] 
                coordinates {(-4,0) (-3,0) (-3,0.166) (3,0.166) (3,0) (4,0)};
                \addlegendentry{\(\mathcal{U}(-3, 3)\)}
        \end{axis}
    \end{tikzpicture}
    \caption{Distribuzioni uniformi}
\end{figure}

Definiamo ora la \textbf{distribuzione gaussiana} (o \textbf{normale}):

\begin{definition}[DISTRIBUZIONE UNIFORME]
    Diciamo che una v.a. $X$ è \textbf{gaussiana} (o \textbf{normale}) se 
    la sua \textbf{densità} è data da:
    \[
        p_X(x) = \frac{1}{\sqrt{2 \pi} \sigma } 
        \exp \left( - \frac{(x-\mu)^2}{2 \sigma^2} \right)
    \] 
\end{definition}

Dove $\sigma= \sqrt{\sigma_x^2}$ è la \textbf{deviazione standard} della
distribuzione e $\mu = m_x$ è la sua \textbf{media}.

Similmente alla notazione usata sopra, scriviamo $x\sim \mathcal{N}(\mu , \sigma)$ 
per indicare che $x$ è campionata da una distribuzione uniforme con 
media $\mu $ e deviazione standard $\sigma$ (e quindi varianza $\sigma^2$).

\begin{figure}[!ht]
    \centering
    \begin{tikzpicture}
        \begin{axis}[
            width=0.8\textwidth,
            height=6cm,
            axis lines=middle,
            xlabel={$x$},
            ylabel={$p_X(x)$},
            xmin=-6, xmax=6,
            ymin=0, ymax=1.2,
            xtick={-6,-4,-2,0,2,4,6},
            ytick={0, 0.2, 0.4, 0.6, 0.8, 1.0},
            grid=major,
            grid style={dashed, gray!30},
            xlabel style={below right},
            ylabel style={above left},
            legend pos=north east,
            legend style={font=\small, draw=none, fill=none},
            % Definiamo una funzione per la gaussiana per rendere il codice più pulito
            declare function={
                gauss(\x,\mu,\sigma) = (1/(\sigma*sqrt(2*pi))) * exp(-((\x-\mu)^2)/(2*\sigma^2));
            }
            ]
            % Gaussiana Standard (mu=0, sigma=1)
            \addplot[thick, nord10, domain=-6:6, samples=200] {gauss(x, 0, 1)};
                \addlegendentry{\(\mathcal{N}(0, 1)\)}

            % Gaussiana stretta e spostata (mu=-2, sigma=0.6)
            \addplot[thick, nord11, domain=-6:6, samples=200] {gauss(x, -2, 0.6)};
                \addlegendentry{\(\mathcal{N}(-2, 0.6)\)}

            % Gaussiana larga e spostata (mu=1.5, sigma=1.8)
            \addplot[thick, nord14, domain=-6:6, samples=200] {gauss(x, 1.5, 1.8)};
                \addlegendentry{\(\mathcal{N}(1.5, 1.8)\)}
        \end{axis}
    \end{tikzpicture}
    \caption{Distribuzioni gaussiane}
\end{figure}

Molto spesso è utile considerare la cosiddetta \textbf{normale standard}
(o gaussiana standard):

\begin{definition}[NORMALE STANDARD]
    Una \textbf{distribuzione normale} con \textbf{media nulla} e
    \textbf{varianza unitaria} si dice \textbf{normale standard}:
    \[
        \mathcal{N}(0,1)
    \] 
\end{definition}

Per distribuzioni gaussiane standard, i valori della funzione
di \textbf{distribuzione cumulativa} $\Phi(a)$ sono \textbf{tabulati}.
\[
    \Phi(a) = P[X \le a] = \int_{-\infty}^a \frac{1}{\sqrt{2\pi}} e^{-0.5 x^2} dx 
\]

Spesso risulta utile considerare il \textbf{valore complementare}.
Anche questi valori sono tabulati e si trovano nelle \textit{tabelle della funzione Q}:
\[
    Q(a) = 1 - \Phi(a)
\]

\section{Processi aleatori}

I processi aleatori possono essere interpretati come 
\textbf{un'estensione del concetto di variabile aleatoria} 
quando viene \textbf{preso in considerazione anche il parametro temporale}.

Diamone quindi la definizione:

\begin{definition}[PROCESSO ALEATORI]
    Si definisce \textbf{processo aleatorio} $v(t)$ (o \textit{stocastico}) una
    \textbf{famiglia di variabili aleatorie} dipendenti dal tempo:
    \[
        v(t) \coloneqq  \{X(t), t\in T \subseteq \mathbb{R}_+\}  
    \] 
\end{definition}

Ciascuna variabile è definita su uno \textit{spazio campionario} $\Omega$ e
\textit{assume valori} in un insieme definito \textbf{spazio degli stati del processo}.

Lo spazio della variabile \textbf{tempo} può essere \textbf{continuo o discreto}.
Nel primo caso si parla di \textit{processo stocastico tempo-continuo}, mentre
nel secondo di \textit{processo stocastico tempo-discreto}.

\begin{remark}
    Notiamo che, \textbf{fissato un istante} $t$, $v(t)$
    \textit{"decade"} ad una semplice \textbf{variabile aleatoria}:
    un processo aleatorio \textbf{associa} quindi una
    \textbf{variabile aleatoria} a \textbf{ciascun istante}
    temporale.
\end{remark}

\subsection{Distribuzione, densità di probabilità e media}

In modo analogo a quanto fatto per le variabili aleatorie, si
definiscono:

\begin{definition}[DISTRIBUZIONE DI PROBABILITA' DEL PRIMO ORDINE]
    \[
        F_v (a,t) \coloneqq P[v(t) \le a]
    \] 
\end{definition}

\begin{definition}[DENSITA' DI PROBABILITA' DEL PRIMO ORDINE]
    \[
        p_v(a,t) \coloneqq \frac{d}{da} F_v(a,t)
    \] 
\end{definition}

Chiaramente, vale sempre la condizione di \textbf{normalizzazione}:
\[
    \int_{-\infty}^{+\infty} p_v (a,t) da = 1
\]

Sempre in modo analogo a quanto visto per le v.a., definiamo la
\textbf{media} di un processo aleatorio:

\begin{definition}[MEDIA DI UN PROCESSO ALEATORIO]
    \[
        m_v(t) \coloneqq \mathbb{E}[v(t)] = \int_{-\infty}^{+\infty} a p_v(a,t) da
    \] 
\end{definition}

Possiamo dare la seguente interpretazione alla media di un processo aleatorio:

\begin{remark}
    La \textbf{media} di un processo aleatorio è una \textbf{funzione}
    che associa a \textbf{ciascun istante} il \textbf{valore atteso del processo}.
\end{remark}

Per concludere, diamo la definzione di \textbf{processi stazionari}:

\begin{definition}[PROCESSO STAZIONARIO]
    Un processo aleatorio si dice \textbf{stazionario} se $p_v(a,t)$
    \textbf{non dipende dal tempo} $t$:
    \[
        p_v(a,t) = p_v(a)
    \] 
\end{definition}

Ciò significa che il processo è \textbf{descritto dalla stessa v.a. in ogni istante}.
Chiaramente, ne consegue che la \textbf{media è costante}; infatti:
\[
    m_v(t) = \mathbb{E}[v(t)] = \int_{-\infty}^{+\infty} a p_v(a) da = m_v \text{ , }\forall t
\]

\subsection{Esempio}

Consideriamo il \textbf{processo aleatorio} seguente:
\[
    v(t) = A \cos(2\pi f_{0}t)
\]
In cui $A\sim \mathcal{U}(0,5)$.

\begin{remark}
    Possiamo interpretare $v(t)$ come un \textit{segnale sinusoidale} a
    frequenza $f_{0}$ la cui \textbf{ampiezza varia casualmente} tra 
    $0$ e $5$ (Volt per esempio).
\end{remark}

Poniamo $c\coloneqq \cos (2\pi f_{0}t)$. Abbiamo:
\[
    A \sim \mathcal{U}(0,5) \implies v(t) = Ac \sim \mathcal{U}(0,5c)
\]
Per cui abbiamo:
\[
    p_v(t) = \begin{cases}
        \frac{1}{5c} & 0 \le a \le 5c \\ 
        0 & \text{altrimenti}
    \end{cases}
\]
Che possiamo riscrivere come:
\[
    p_v(t) = \frac{1}{5 \cos(2\pi f_{0}t)} \text{rect} \left( 
               \frac{a - 2.5 \cos(2\pi f_{0}t)}{5 \cos(2\pi f_{0}t)} \right)
\]
Calcoliamo $p_v(t)$ per alcuni valori di $t$, per esempio:
\[
    \begin{aligned}
        p_v(a,0) &= \frac{1}{5} \text{rect} \left( \frac{a-2.5}{5} \right) \\ 
        p_v \left(a, \frac{1}{4f_{0}} \right) &= \delta(a)
    \end{aligned}
\]

La \textbf{media} risulta invece:
\[
    m_v(t) = \mathbb{E}[v(t)] = \int_{-\infty}^{+\infty} a p_v(a,t) da = 2.5 \cos(2\pi f_{0}t)
\]
Che è un \textit{risultato intuitivo}, dato che l'ampiezza $A$ \textit{varia uniformemente}.
